% ============================================================
%  JCIM Submission — LaTeX Source
%  Title: Leakage-Controlled Fragment Replacement Ranking
%         with Audited Candidate-Level Scoring
%
%  Generated: 2026-06-02 by PaperSpine LaTeX
%  Template: Standard article (adaptable to achemso)
% ============================================================

\documentclass[12pt,letterpaper]{article}

% --- Packages ---
\usepackage[utf8]{inputenc}
\usepackage[T1]{fontenc}
\usepackage{amsmath,amssymb}
\usepackage{graphicx}
\usepackage{booktabs}
\usepackage[margin=1in]{geometry}
\usepackage{hyperref}
\usepackage[numbers,sort&compress]{natbib}
\usepackage{xspace}
\usepackage{multirow}
\usepackage{array}

% --- Metadata ---
\title{Leakage-Controlled Fragment Replacement Ranking with Audited Candidate-Level Scoring}

\date{}

\begin{document}
\maketitle

% ============================================================
%  ABSTRACT
% ============================================================

\begin{abstract}
Matched molecular pair (MMP)-derived fragment replacement ranking is useful for studying observed substitution patterns, but random splits can leak fragment-attachment transforms across evaluation sets, and learned models can exploit dataset-specific shortcut features that fail under distribution shift. We present a transform-heldout benchmark for closed-vocabulary scaffold-conditioned replacement ranking with 13,347 secondary blind queries, together with a candidate-level histogram gradient-boosted scorer built from base-ranker outputs, train-derived priors, molecular descriptors, and frozen categorical features. A post-audit 77-feature scorer achieves Top-10 = 0.9243 (95\% CI [0.9199, 0.9287]), improving over the Score Blend baseline by $\Delta$Top-10 = 0.0686 (95\% CI [0.0638, 0.0733]).

The central methodological finding is a deletion: removing five per-query prior-rank features eliminates a non-transferable shortcut, improving blind Top-10 by +0.0393 (95\% CI [0.0354, 0.0432]) over the initial 82-feature scorer and reducing Score Blend hit loss from 596 to 101 queries (5.9$\times$ reduction). Because this deletion was identified after blind diagnostics, the 77-feature result is a locked post-selection finding that motivates prospective replication rather than a fully prospective feature-selection result. A dual-mode A4C computational review proxy annotates provenance strata and flags structural alerts as unvalidated computational triage signals. All labels are structure-derived from ChEMBL MMPs and do not establish activity preservation.
\end{abstract}

% ============================================================
%  1. INTRODUCTION
% ============================================================

\section{Introduction}

Large chemical databases enable the extraction of many structure-derived substitution patterns from matched molecular pair analysis, yet computational methods for ranking fragment replacements lack evaluation protocols that prevent models from exploiting dataset-specific regularities rather than learning transferable replacement patterns. We study closed-vocabulary scaffold-conditioned fragment replacement proposal: given an old fragment and its attachment signature, a model ranks candidate replacements from the attachment-compatible subset of a global training replacement vocabulary. The task is deliberately structure-derived---it asks whether known replacement patterns can be recovered under controlled evaluation, not whether a proposed replacement preserves activity.

Building a credible benchmark is difficult for two main reasons. First, random data splits can leak the same fragment-attachment transform into both training and evaluation sets, allowing models to succeed through memorization rather than generalization. Prior work has demonstrated that random splits in ligand-based classification tasks can substantially overestimate model performance~\cite{wallach2018most}, and systematic data leakage is a recognized cause of reproducibility failures in machine-learning-based science~\cite{kapoor2023leakage}. Second, even under a leakage-controlled split, learned models can exploit dataset-specific shortcut features---those that are predictive within the development distribution but whose meaning changes under fragment identity shift---and these shortcuts are difficult to identify without systematic post-audit ablation. Shortcut learning, where models exploit spurious correlations that fail to transfer across distributions, is a well-documented failure mode across deep learning~\cite{geirhos2020shortcut}.

We address these issues with two contributions. First, we introduce a leakage-controlled benchmark for closed-vocabulary scaffold-conditioned fragment replacement ranking. The benchmark uses a transform-heldout split that isolates unseen fragment-attachment combinations, preventing memorization of transform identities. After post-audit feature pruning, a 77-feature HistGB scorer achieves Top-10 = 0.9243 on the secondary blind evaluation protocol (13,347 queries). Second, through systematic post-audit ablation, we identify five per-query prior-rank features as non-transferable shortcuts: removing them improves blind Top-10 by +0.0393 over the initial 82-feature scorer, reduces baseline hit loss by 5.9$\times$, and reveals a design rule---sparse prior-score rank transforms can become generalization hazards under fragment distribution shift. We enforce strict protocol separation between prospective evaluation, post-audit analysis, and workflow interpretation so that each evidentiary role is explicit. A dual-mode A4C computational review proxy provides provenance-stratified alert reporting as a supporting workflow feature, with the explicit caveat that alert rates are unvalidated computational screening signals.

On the secondary blind evaluation protocol, the post-audit 77-feature scorer reaches Top-10 = 0.9243 (95\% CI [0.9199, 0.9287]), improving over the Score Blend baseline by $\Delta$Top-10 = 0.0686 (95\% CI [0.0638, 0.0733]). Removing the five prior\_ranks features yields a paired gain of $\Delta$Top-10 = 0.0393 (95\% CI [0.0354, 0.0432]) over the 82-feature scorer. Because these features were deleted after blind diagnostics, we treat the 77-feature model as a post-audit locked result and the prior-rank interpretation as an analytical finding. Under the A4C proxy, G2 and G3 provenance groups are fully covered with alert rates of 46.85\% and 9.67\%, respectively; G4 has sparse A4C coverage (5.63\%), so no reliable G4-wide alert estimate is possible.

Our contributions are as follows. First, we provide a leakage-controlled benchmark for MMP-derived scaffold-conditioned fragment replacement with a transform-heldout split and secondary blind evaluation protocol. Second, we present a candidate-level scorer that improves over base-ranker fusion and show, through a post-audit ablation with paired bootstrap uncertainty, that sparse prior-rank features can become non-transferable shortcuts---the most impactful feature-engineering decision being a deletion rather than an addition. Third, we enforce strict protocol separation: prospective performance evidence, post-audit mechanism analysis, and workflow interpretation each occupy distinct evidentiary roles. Fourth, we frame exploratory outputs through a dual-mode workflow that keeps provenance and alert stratification explicit while acknowledging the computational proxy's limitations.

% ============================================================
%  2. RELATED WORK
% ============================================================

\section{Related Work}

\textbf{Matched molecular pairs and structure-derived replacements.}
Bioisosteric replacement is a longstanding design concept in medicinal chemistry~\cite{patani1996bioisosterism}. Matched molecular pair analysis extracts local substitution patterns from large compound collections by identifying molecules that differ at a single structural site while sharing an identical scaffold~\cite{hussain2010computationally,griffen2011matched}. Databases such as ChEMBL~\cite{zdrazil2024chembl} make structure-derived replacement catalogues feasible at scale. These data record observed structural substitutions and serve as supervision for replacement proposal, but they do not encode activity continuity---a limitation inherited by any model trained on them.

\textbf{Rule-based and database-driven tools.}
SwissBioisostere organizes observed replacements into a searchable database for ligand design~\cite{wirth2013swissbioisostere}, while CReM uses fragment environments from known compounds to perform context-aware replacement generation~\cite{polishchuk2020crem}. NeBULA extends this paradigm by mining over 700 medicinal chemistry references to compile bioisosteric replacement reaction rules encoded as SMARTS and SMIRKS patterns, covering multiple replacement categories from scaffold hopping to general isosteric substitution~\cite{huang2025nebula}. These frameworks exploit empirical replacement context at scale, but they define different task formulations---typically open-ended generation or database lookup---rather than the closed-vocabulary query-level ranking benchmark studied here.

\textbf{Learning-based methods.}
Deep learning methods for bioisosteric replacement include descriptor-based neural networks~\cite{ertl2020identification} and deep generative models that autonomously select fragments for removal and insertion while controlling multiple properties such as logP, drug-likeness, and synthetic accessibility simultaneously~\cite{kim2025deepbioisostere}. GraphBioisostere applies a Siamese graph neural network with attention-based global pooling (MMGX encoder) to ChEMBL-derived molecular pairs, training on 779,081 compound pairs across 1,897 targets, and demonstrates transfer learning for target-specific potency change prediction across five pharmaceutically relevant targets~\cite{masunaga2026graphbioisostere}. These methods show that structural representations capture replacement-relevant information beyond frequency statistics, but they are typically evaluated on random splits that may overestimate generalization. Our work complements the generative formulation with a ranking perspective that explicitly separates candidate scoring from candidate generation.

\textbf{Rank aggregation and fusion.}
Borda count is a classical parameter-free aggregation rule for combining ranked lists~\cite{dwork2001rank}. Reciprocal rank fusion~\cite{cormack2009reciprocal} provides a weighted alternative. Our use of fusion is narrower: we test whether structural and empirical rankers remain complementary under leakage-controlled evaluation and use Borda's parameter-free property as a practical necessity---the transform-heldout split discourages development-set-based tuning of fusion weights.

\textbf{Computational alert and risk filters.}
PAINS~\cite{baell2010new} and Brenk alerts~\cite{brenk2008lessons} are widely used medicinal-chemistry screening heuristics. Recent systematic activity-profile studies demonstrate that bioisosteric replacements, even structurally conservative ones, can produce unexpected off-target effects~\cite{helmke2025data}. We adopt a similar philosophy for the A4C proxy: a rule-based annotation layer that organizes provenance strata and flags candidate-level structural alerts without converting exploratory outputs into approval decisions.

% ============================================================
%  3. METHODS
% ============================================================

\section{Methods}

We present the complete methodological framework. Section~3.1 formalizes the ranking task, supervision labels, and evaluation metrics. Section~3.2 describes benchmark construction. Sections~3.3--3.5 cover the ranking methods: five base rankers (Section~3.3), a candidate-level scorer with 77 features and a gradient boosting architecture (Section~3.4), and a leave-one-family-out ablation that identifies a feature family harmful to cross-split generalization (Section~3.5). Sections~3.6--3.7 describe the evaluation protocol and the A4C computational review proxy.

\subsection{Task Formulation}

We formalize scaffold-conditioned fragment replacement as a closed-vocabulary, multi-positive learning-to-rank problem. The candidate vocabulary $\mathcal{V}_{\mathrm{train}}$ is training-derived and closed: it is constructed exclusively from the training split. Transform identities $(f^{\mathrm{old}}, \sigma)$ are held out across splits (Section~3.2). Let $\mathcal{Q} = \{q_i\}_{i=1}^{N}$ denote a query set. Each query pairs an old fragment with its attachment context,
%
\begin{equation}
q_i = \bigl(f_i^{\mathrm{old}}, \sigma_i\bigr),
\end{equation}
%
where $f_i^{\mathrm{old}}$ is the fragment to be replaced and $\sigma_i$ is the attachment signature---the atom and bond context through which the fragment is connected to the scaffold. For query $q_i$, only candidates that are attachment-compatible with $\sigma_i$ are eligible:
%
\begin{equation}
\mathcal{C}_i = \bigl\{c\in\mathcal{V}_{\mathrm{train}}:\chi(c,\sigma_i)=1\bigr\},
\end{equation}
%
where $\chi(c,\sigma_i)\in\{0,1\}$ is a deterministic attachment-compatibility predicate encoding valence and bond-order constraints.

\textbf{Supervision labels.}
Let $\mathcal{M}$ denote the set of structure-derived replacement triples from MMP extraction after the train/evaluation split has been fixed. For query $q_i$, the positive set is $\mathcal{P}_i = \{c\in\mathcal{C}_i: (f_i^{\mathrm{old}},\sigma_i,c)\in\mathcal{M}\}$, with candidate-level binary label $y_{ic} = \mathbb{I}[c\in\mathcal{P}_i]$. Queries can be multi-positive. Labels are structure-derived: they indicate observed structural substitution, not activity preservation.

\textbf{Evaluation metrics.}
The best positive rank for query $q_i$ under model $\theta$ is $r_i^+(\theta)=\min_{c\in\mathcal{P}_i}\rho_\theta(i,c)$. The per-query hit at cutoff $K$ is $h_i^{(K)}(\theta)=\mathbb{I}[r_i^+(\theta)\le K]$, with a hit scored if any positive fragment appears in the top $K$. Query-level Top-$K$ accuracy and mean reciprocal rank are $\operatorname{Top@K}(\theta;\mathcal{Q}) = \frac{1}{N}\sum_{i=1}^{N}h_i^{(K)}(\theta)$ and $\operatorname{MRR}(\theta;\mathcal{Q}) = \frac{1}{N}\sum_{i=1}^{N}1/r_i^+(\theta)$. Top-10 is the primary metric; Top-1, Top-5, Top-20, Top-50, and MRR are secondary diagnostics. Confidence intervals use nonparametric bootstrap at the query level (Section~3.6.2).

\textbf{Scoring model and induced rank.}
A scoring model with parameters $\theta$ assigns each candidate a real-valued score $s_{\theta} : (q_i,c) \mapsto s_{\theta}(i,c)\in\mathbb{R}, c\in\mathcal{C}_i$. Candidates are ranked in decreasing order of $s_{\theta}(i,c)$. To make the ranking function well-defined under score ties, let $\tau_i(c)$ be a fixed deterministic tie-breaking key. The induced rank is
%
\begin{equation}
\rho_{\theta}(i,c) = 1 + \sum_{c'\in\mathcal{C}_i} \mathbb{I}\!\bigl[ s_{\theta}(i,c')>s_{\theta}(i,c) \;\lor\; \bigl(s_{\theta}(i,c')=s_{\theta}(i,c)\land \tau_i(c')<\tau_i(c)\bigr) \bigr].
\end{equation}

\subsection{Benchmark Construction}

\textbf{Data source and MMP extraction.}
We use ChEMBL37K~\cite{zdrazil2024chembl}, a curated subset of approximately 37,000 drug-like molecules from ChEMBL33. MMPs are extracted by applying retrosynthetic bond-cleavage rules~\cite{hussain2010computationally,griffen2011matched}: each molecule is recursively decomposed at acyclic single bonds to generate fragment--scaffold pairs; any two molecules sharing identical scaffold fragments but bearing different substituent fragments at the same attachment point are recorded as an MMP. Quality filters include salt stripping, stereochemistry normalization, removal of molecules with PAINS or Brenk alerts~\cite{baell2010new,brenk2008lessons}, and exclusion of fragments with molecular weight below 15~Da or above 250~Da.

\textbf{Decoy construction.}
Negatives are constructed through a decoy repair procedure: for each query, $f_i^{\mathrm{old}}$ is paired with a fragment $c_{\mathrm{decoy}} \in \mathcal{V}_{\mathrm{train}}$ attached to a non-cognate scaffold. Four stages: (1) 1:1 positive-to-decoy ratio; (2) wrong-positive removal; (3) deduplication; (4) a 1:5 diagnostic set for class-imbalance stress testing.

\textbf{Split design: transform-heldout.}
A conventional random split would allow the same $(f^{\mathrm{old}}, \sigma)$ pair---the \emph{transform identity}---to appear in both training and test sets. We prevent this with a \textbf{transform-heldout split}: all unique $(f^{\mathrm{old}}, \sigma)$ pairs are partitioned such that no pair appears in both training and evaluation splits. We construct three evaluation tiers: (i) \textbf{Development/calibration set}---a held-out subset of the training partition for architecture, feature, and hyperparameter selection; (ii) \textbf{Secondary blind set}---13,347 queries, zero $(f^{\mathrm{old}}, \sigma)$ overlap, primary evaluation benchmark; (iii) \textbf{Canonical analysis set}---21,052 queries for robustness and mechanism analysis only.

\begin{table}[ht]
\centering
\caption{Candidate-matrix statistics for the closed-vocabulary ranking benchmark.}
\label{tab:M1}
\begin{tabular}{lrrrrrrr}
\toprule
Split & Rows & Queries & Pos. rows & Old frags & Transforms & Single-pos & Multi-pos \\
\midrule
Train & 16,497,750 & 109,985 & 181,483 & 109 & 351 & 73,870 & 36,115 \\
Dev./calibration & 2,006,250 & 13,375 & 21,408 & 8 & 27 & 8,686 & 4,689 \\
Secondary blind & 2,002,050 & 13,347 & 20,082 & 19 & 63 & 9,377 & 3,970 \\
\bottomrule
\end{tabular}
\end{table}

\begin{table}[ht]
\centering
\caption{Leakage verification: overlap counts of $(f^{\mathrm{old}}, \sigma)$ transform identities.}
\label{tab:M2}
\begin{tabular}{lrr}
\toprule
Split Pair & Transform Overlap & Verified \\
\midrule
Train / Development set & 0 & Yes \\
Train / Secondary blind & 0 & Yes \\
Train / Canonical test & 0 & Yes \\
Canonical / Secondary blind queries & 0 & Yes \\
\bottomrule
\end{tabular}
\end{table}

\subsection{Base Ranking Methods}

We implement five baseline ranking methods spanning a range of inductive biases. Each captures a distinct signal: global replacement statistics, learned structural compatibility, interpretable molecular properties, and their consensus. During development, query-level routing strategies that select among base rankers on a per-query basis were explored but did not improve over per-candidate scoring; the consistent failure of query-level routing to outperform candidate-level scoring is analyzed in Section~5.2.

\subsubsection{Attachment-Frequency Ranker}

The Attachment-Frequency ranker estimates the empirical training-set prior. Let $n_{\mathrm{train}}(c,\sigma)$ be the count of training replacement pairs where $c$ appears under attachment $\sigma$:
%
\begin{equation}
s_{\mathrm{att}}(c,\sigma) = \widehat{p}_{\mathrm{att}}(c\mid\sigma) = \frac{n_{\mathrm{train}}(c,\sigma)}{\sum_{c'\in\mathcal{V}_{\mathrm{train}}} n_{\mathrm{train}}(c',\sigma)}.
\end{equation}
%
The method has no learned parameters and serves as a simple empirical frequency baseline.

\subsubsection{Dual Encoder (DE) Ranker}

The DE ranker learns a measure of structural compatibility through a dual-encoder architecture. The query encoder takes the Morgan fingerprint (ECFP, radius 2, 2048 bits) of $f_i^{\mathrm{old}}$ and a learned embedding of $\sigma_i$, concatenates them, and passes the result through a two-layer MLP (hidden dimensionality $d=128$, ReLU). The candidate encoder independently maps the candidate's Morgan fingerprint through a separate two-layer MLP (same architecture, not weight-tied):
%
\begin{equation}
\mathbf{h}^{q}_{i} = g_{\phi}\!\left(x(f_i^{\mathrm{old}}),a(\sigma_i)\right), \qquad \mathbf{h}^{c}_{c} = h_{\psi}\!\left(x(c)\right).
\end{equation}
%
The DE score is cosine similarity: $s_{\mathrm{DE}}(i,c) = \langle \mathbf{h}^{q}_{i},\mathbf{h}^{c}_{c}\rangle / (\lVert\mathbf{h}^{q}_{i}\rVert_2\,\lVert\mathbf{h}^{c}_{c}\rVert_2)$. Training uses a margin-based ranking loss with $K=20$ negatives per positive and margin $\delta=0.3$:
%
\begin{equation}
\mathcal{L}_{\mathrm{DE}} = \sum_{i=1}^{N} \sum_{c^{+}\in\mathcal{P}_i} \sum_{c^{-}\in\mathcal{N}_i} \left[\delta - s_{\mathrm{DE}}(i,c^{+}) + s_{\mathrm{DE}}(i,c^{-})\right]_{+}.
\end{equation}

\subsubsection{Histogram Gradient-Boosted Ranker (HGB)}

The HGB ranker uses explicit molecular features organized into four groups: (1) frequency features; (2) attachment signature features; (3) molecular property descriptors (heavy atom count, molecular weight, logP, TPSA, ring count, HBD/HBA counts, rotatable bonds---for both candidate and query fragments); (4) fingerprint similarity (Tanimoto coefficient between Morgan fingerprints, radius 2, 2048 bits). The model is \texttt{HistGradientBoostingClassifier} (scikit-learn) with 200 iterations, max depth 6, learning rate 0.1. Candidates are ranked by predicted probability.

\subsubsection{Borda Fusion and Score Blend}

\textbf{Borda fusion.} For query $q_i$, let $\rho_m(i,c)$ be the rank assigned by method $m\in\{\mathrm{DE},\mathrm{HGB}\}$:
%
\begin{equation}
S_{\mathrm{Borda}}(i,c) = \sum_{m\in\{\mathrm{DE},\mathrm{HGB}\}} \left(|\mathcal{C}_i| + 1 - \rho_m(i,c)\right).
\end{equation}
%
The parameter-free design is deliberate: the transform-heldout split discourages development-set-based tuning of fusion weights.

\textbf{Score Blend (MLP + HGB).} A rank-only MLP takes three per-candidate rank values (DE, HGB, Attachment-Frequency) and outputs $s_{\mathrm{MLP}}(q_i, c)$. The blended score is $s_{\mathrm{blend}}(i,c) = \lambda\, z_i\{s_{\mathrm{MLP}}\}(c) + (1-\lambda)\, z_i\{s_{\mathrm{HGB\text{-}refit}}\}(c)$, with $\lambda=0.95$.

\subsection{Candidate-Level Scorer}

\subsubsection{Architecture}

The base rankers provide broad but shallow assessments. We introduce a \textbf{candidate-level scorer} that combines 77 features with histogram gradient boosting. Given query $q_i$ and candidate $c$, the scorer computes $\boldsymbol{\phi}_{ic} = \Phi(q_i,c) \in \mathbb{R}^{77}$ and maps it to a probability: $\widehat{p}_{ic} = F_{\Theta}(\boldsymbol{\phi}_{ic})$. The training target is $y_{ic}$, fit by weighted binary log loss. We choose HistGB for three reasons: (1) mixed-type inputs are handled natively; (2) non-linear feature interactions are captured; (3) built-in feature importance estimates support ablation.

\subsubsection{Feature Families}

The 77 features organize into 9 families (Table~\ref{tab:M3}) plus 22 one-hot categorical features. Five prior\_ranks features were removed post-audit as non-transferable shortcuts (Section~3.5).

\begin{table}[ht]
\centering
\caption{Feature families of the candidate-level scorer with audit trail.}
\label{tab:M3}
\small
\begin{tabular}{lrlcccc}
\toprule
Family & Count & Source & Train-Only & Label-Free & Retained & Leakage Control \\
\midrule
F1: Model Scores & 5 & Base ranker outputs ($z$-normalized) & Yes & Yes & Yes & Train-set statistics \\
F2: Model Raw Scores & 4 & Base ranker raw outputs & Yes & Yes & Yes & No eval-set dependence \\
F3: Model Ranks & 6 & Per-query ranks from base rankers & Yes & Yes & Yes & Query-relative \\
F4: Top-K Flags & 4 & Binary: candidate in base ranker top-10 & Yes & Yes & Yes & Dev-set thresholds \\
F5: Prior Scores & 5 & Logit-transformed empirical priors & Yes & Yes & Yes & Train counts only \\
F6: Prior Statistics & 12 & Smoothed rates, support/pos. counts & Yes & Yes & Yes & Train counts \\
F7: PMI / Contrast & 6 & PMI and contrast across prior levels & Yes & Yes & Yes & Train co-occurrence \\
F8: Mol. Descriptors & 10 & RDKit: candidate + deltas from query & No & Yes & Yes & Structure-derived \\
F9: Similarity \& Freq & 3 & Tanimoto, global/attach frequency & Yes & Yes & Yes & Train counts \\
Categorical (one-hot) & 22 & Fragment ID, attach sig., freq bin & Yes & Yes & Yes & Frozen schema \\
Prior\_Ranks (removed) & 5 & Per-query rank of prior scores & Yes & Yes & No & Shortcut hazard \\
\midrule
Total Retained & 77 & & & & & \\
\bottomrule
\end{tabular}
\end{table}

\subsubsection{Training Protocol}

The scorer uses \texttt{HistGradientBoostingClassifier} (scikit-learn~\cite{pedregosa2011scikit}) with max 200 iterations, max depth 6, learning rate 0.1, \texttt{class\_weight='balanced'}, \texttt{random\_state=42}. Training uses 5-fold GroupKFold grouped by query identifier. Negative subsampling at 5:1. All hyperparameters were fixed on the development set.

\subsubsection{Frozen Categorical Schema Alignment}

Categorical features pose a schema-mismatch risk because development and blind sets can contain disjoint fragment identities. We use a frozen categorical schema: the one-hot encoder is fitted on the development set; unseen categories map to OTHER; output dimensionality is forced to match the development schema.

\subsection{Feature Ablation: The Prior\_Ranks Removal}

We used leave-one-family-out ablation to identify the post-audit feature-pruned configuration. The base ranker architectures, scorer feature families F1--F9, and all training hyperparameters were configured on the development set. The prior\_ranks removal was identified during post-audit analysis of the initial 82-feature scorer after blind-set diagnostics revealed fragment-specific degradation. We therefore report the 77-feature scorer as a post-audit locked model, not as a fully prospective pre-registered feature-selection result.

\subsubsection{Ablation Design}

During initial development, five additional features encoding per-query rank positions of prior scores were added to the base 77-feature set, yielding an 82-feature configuration. These five features, collectively denoted \textbf{prior\_ranks}, are backoff\_logit\_rank, cont\_prior\_rank, t\_p4\_logit\_rank, p1\_logit\_rank, and p3\_logit\_rank. For prior score $u_j(i,c)$, the rank feature is:
%
\begin{equation}
r_j^{\mathrm{prior}}(i,c) = 1 + \sum_{c'\in\mathcal{C}_i} \mathbb{I}\!\left[ u_j(i,c')>u_j(i,c) \;\lor\; \bigl(u_j(i,c')=u_j(i,c)\land \tau_i(c')<\tau_i(c)\bigr) \right].
\end{equation}
%
Let $\mathcal{A}_{82}$ be the index set and $\mathcal{R}_{\mathrm{prior}} \subset \mathcal{A}_{82}$ the five prior\_ranks indices. The 77-feature map is the coordinate projection: $\boldsymbol{\phi}^{(77)}_{ic} = \Pi_{\mathcal{A}_{82}\setminus\mathcal{R}_{\mathrm{prior}}}(\boldsymbol{\phi}^{(82)}_{ic})$.

\subsubsection{Why Prior\_Ranks Degrade Generalization}

The prior\_ranks features differ structurally from all others: they encode per-query rank position, which is \textbf{query-relative}. The development set contains 8 distinct old\_fragment identities; the secondary blind set contains 19. The HistGB model learns thresholds reflecting the ranking patterns of 8 development fragments; these fail to transfer to 19 blind fragments. We interpret this as \textbf{shortcut learning}~\cite{geirhos2020shortcut}: rank features are predictive within the development distribution, so the model defaults to them instead of learning transferable chemical patterns.

\textbf{Design rule.} Sparse prior-score rank transforms can become generalization hazards in cross-split molecular ranking; raw scores, rates, and intrinsic molecular properties are safer default representations under fragment distribution shift.

\subsection{Evaluation Protocol}

\subsubsection{Metrics and Ranking Evaluation}

For each held-out query, candidates are sorted by descending $s_\theta(i,c)$. Multi-positive queries score a hit if any positive candidate appears in the top $K$. Primary metric: Top-10 accuracy. Secondary: Top-1, Top-5, Top-20, Top-50, MRR.

\subsubsection{Bootstrap Confidence Intervals}

95\% confidence intervals use nonparametric bootstrap with $B=5{,}000$ replicates, resampling at the query level. For the $b$-th replicate, $\mathcal{Q}^*_b$ of size $N$ is drawn with replacement from $\mathcal{Q}$: $\hat{\theta}^*_b = \frac{1}{N}\sum_{i\in\mathcal{Q}^*_b} h_i^{(K)}(\theta)$, $\operatorname{CI}_{95\%}(\theta) = [\hat{\theta}^*_{(0.025)}, \hat{\theta}^*_{(0.975)}]$. For pairwise comparisons, a paired difference bootstrap is used: $\hat{\Delta}^*_b = \hat{\theta}^*_b(m) - \hat{\theta}^*_b(b)$.

\subsubsection{Rescue and Lost Analysis}

Let $b$ be a baseline and $m$ the model under evaluation:
%
\begin{align}
\operatorname{Rescue}_i(m,b) &= \mathbb{I}\!\left[r_i^+(b)>10\right]\; \mathbb{I}\!\left[r_i^+(m)\le 10\right], \\
\operatorname{Lost}_i(m,b) &= \mathbb{I}\!\left[r_i^+(b)\le 10\right]\; \mathbb{I}\!\left[r_i^+(m)>10\right].
\end{align}
%
Aggregated counts $R(m,b)$, $L(m,b)$, and net gain $G=R-L$ decompose accuracy changes into recoveries and regressions.

\subsubsection{Protocol Separation}

We enforce strict separation of evaluation protocols by evidentiary role (Table~\ref{tab:M4}). Primary claims cite only secondary blind results.

\begin{table}[ht]
\centering
\caption{Evaluation protocol separation.}
\label{tab:M4}
\begin{tabular}{lrll}
\toprule
Protocol & Queries & Purpose \\
\midrule
Development/calibration & -- & Architecture, feature, hyperparameter selection \\
Secondary blind & 13,347 & \textbf{Primary performance claim} \\
Canonical analysis & 21,052 & Robustness and mechanism analysis only \\
Dual-mode risk-stratification & -- & Workflow and risk interpretation only \\
\bottomrule
\end{tabular}
\end{table}

\subsubsection{Leakage Control}

All statistics (prior probabilities, smoothing parameters, PMI, $z$-score statistics) derive exclusively from the training split. The frozen categorical schema ensures consistent dimensionality. Zero $(f^{\mathrm{old}}, \sigma)$ overlap between training and any evaluation split (Table~\ref{tab:M2}). All 77 features are computable without evaluation labels.

\subsection{Computational Review Proxy (A4C)}

The A4C annotation layer applies PAINS alerts~\cite{baell2010new} and Brenk alerts~\cite{brenk2008lessons} to flag structural features. The dual-mode workflow offers Conservative Mode (HGB proposals) and Exploration Mode (Borda proposals). Within Exploration Mode, proposals receive provenance labels (Table~\ref{tab:M5}).

\begin{table}[ht]
\centering
\caption{Provenance groups and definitions.}
\label{tab:M5}
\begin{tabular}{lll}
\toprule
Group & Definition & Interpretation \\
\midrule
G4 & $\mathcal{K}_{\mathrm{HGB}} \cap \mathcal{K}_{\mathrm{Borda}}$ & Shared: frequency--similarity consensus \\
G3 & $\mathcal{K}_{\mathrm{Borda}} \setminus \mathcal{K}_{\mathrm{HGB}}$ (DE-elevated) & Expansion beyond frequency priors \\
G2 & $\mathcal{K}_{\mathrm{Borda}} \setminus (\mathcal{K}_{\mathrm{HGB}} \cup \mathcal{K}_{\mathrm{DE}})$ & Borda-emergent, highest novelty \\
\bottomrule
\end{tabular}
\end{table}

The A4C framework is used exclusively for workflow and risk interpretation. Alert rates are computational screening signals, not experimentally validated determinations. We treat G2/G3/G4 as illustrative provenance strata for downstream inspection, not as validated decision rules.

% ============================================================
%  4. RESULTS
% ============================================================

\section{Results}

We evaluate all methods on the secondary blind protocol (13,347 queries) using query-level Top-10 accuracy as the primary metric, with 95\% bootstrap confidence intervals ($B=5{,}000$).

\subsection{Candidate-Level Scoring Improves Blind Ranking}

The post-audit 77-feature scorer achieves Top-10 = 0.9243, improving over the Score Blend baseline by $\Delta$Top-10 = +0.0686 (95\% paired bootstrap CI [+0.0638, +0.0733]). The initial 82-feature scorer achieves Top-10 = 0.8851 as the prospective result. The candidate-level scoring design was partially motivated by the consistent failure of query-level routing strategies to improve over per-candidate scoring; analysis of this routing gap is deferred to Section~5.2.

\begin{table}[ht]
\centering
\caption{Secondary blind Top-10 accuracy, with 95\% bootstrap confidence intervals.}
\label{tab:1}
\begin{tabular}{lcc}
\toprule
Method & Blind Top-10 & 95\% CI \\
\midrule
Attachment-Frequency & 0.6019 & [0.5933, 0.6104] \\
HGB & 0.7437 & [0.7356, 0.7516] \\
Dual Encoder (DE) & 0.8055 & [0.7986, 0.8122] \\
Borda(DE, HGB) & 0.8384 & [0.8321, 0.8447] \\
MLP (rank-only) & 0.8402 & [0.8339, 0.8466] \\
Score Blend (MLP + HGB) & 0.8558 & [0.8495, 0.8621] \\
Initial 82-feature scorer & 0.8851 & [0.8796, 0.8903] \\
\textbf{Post-audit 77-feature scorer} & \textbf{0.9243} & \textbf{[0.9199, 0.9287]} \\
Best-of-DE+HGB diagnostic & 0.8686 & --- \\
\bottomrule
\end{tabular}
\end{table}

\subsection{Removing Non-Transferable Prior-Rank Shortcuts Improves Blind Transfer}

Removing the five prior\_ranks features improves blind Top-10 from 0.8851 to 0.9243, a gain of +0.0393 (95\% paired bootstrap CI [+0.0354, +0.0432]). No other leave-one-family-out ablation produces a comparable improvement.

\begin{table}[ht]
\centering
\caption{Leave-one-family-out ablation. $\Delta$ is relative to the full 82-feature configuration.}
\label{tab:2}
\begin{tabular}{lrrrl}
\toprule
Condition & Features Removed & Blind Top-10 & $\Delta$ & 95\% CI for $\Delta$ \\
\midrule
Full 82-feature scorer & 0 & 0.8851 & -- & -- \\
\textbf{Drop prior\_ranks} & \textbf{5} & \textbf{0.9243} & \textbf{+0.0393} & \textbf{[+0.0354, +0.0432]} \\
Drop model\_ranks (F3) & 6 & 0.8697 & --0.0154 & not reported \\
Drop model\_scores (F1) & 5 & 0.8610 & --0.0241 & not reported \\
\bottomrule
\end{tabular}
\end{table}

\subsection{Feature Pruning Recovers Baseline Misses While Reducing Regression}

Relative to the Score Blend baseline, the 77-feature scorer rescues 1,016 queries and loses 101, net +915 (rescue-to-loss ratio 10.06, CI [8.25, 12.52]). Relative to the initial 82-feature scorer: rescue 626, lose 102, net +524. The 82-feature model had 596 lost queries against Score Blend; the 77-feature scorer reduces this to 101---a 5.9$\times$ reduction.

\begin{table}[ht]
\centering
\caption{Rescue/lost analysis for the post-audit 77-feature scorer. $N = 13{,}347$.}
\label{tab:3}
\begin{tabular}{lrrrrr}
\toprule
Reference & Ref Top-10 & Rescue & Lost & Net & Arithmetic \\
\midrule
Score Blend & 0.8558 & 1,016 & 101 & +915 & Yes \\
Initial 82-feature scorer & 0.8851 & 626 & 102 & +524 & Yes \\
Borda (candidate-matrix) & 0.8456 & 1,152 & 101 & +1,051 & Yes \\
HGB-refit (candidate-matrix) & 0.8623 & 977 & 150 & +827 & Yes \\
DE & 0.8055 & 1,754 & 168 & +1,586 & Yes \\
\bottomrule
\end{tabular}
\end{table}

\subsection{Categorical Schema Alignment Prevents Invalid Blind Prediction}

A silent prediction bug from naive one-hot encoding produced a blind Top-10 of 0.7120. The frozen categorical schema corrected this to 0.8851. This is an implementation audit finding, not a model contribution.

\subsection{Dual-Mode Provenance Strata Provide Coverage-Limited Triage Signals}

\begin{table}[ht]
\centering
\caption{A4C alert rates by provenance group.}
\label{tab:4}
\begin{tabular}{llcc}
\toprule
Group & Definition & A4C Coverage & Alert Rate \\
\midrule
G4 & Shared (HGB $\cap$ Borda) & 5.63\% & 17.60\% among covered \\
G3 & DE-elevated (Borda $\setminus$ HGB) & 100\% & 9.67\% \\
G2 & Borda-emergent (Borda $\setminus$ (HGB $\cup$ DE)) & 100\% & 46.85\% \\
\bottomrule
\end{tabular}
\end{table}

G2 and G3 have complete A4C coverage. G4 has sparse coverage (5.63\%); 94.37\% of G4 candidates have unknown A4C status, so no reliable G4-wide alert estimate is possible.

% ============================================================
%  5. DISCUSSION
% ============================================================

\section{Discussion}

\subsection{Principal Findings}

This study makes three central empirical contributions. First, under a leakage-controlled transform-heldout benchmark, the post-audit 77-feature scorer achieves Top-10 = 0.9243, improving over Score Blend by $\Delta$Top-10 = +0.0686 [95\% CI: +0.0638, +0.0733]. Second, removing five per-query prior-rank features improves blind Top-10 by +0.0393 [95\% CI: +0.0354, +0.0432], traced to shortcut learning. Third, the dual-mode workflow provides coverage-limited computational triage (G2: 46.85\%, G3: 9.67\%).

The prior\_ranks removal was identified post-audit; the 77-feature scorer is a locked post-selection result. The mechanistic explanation---shortcut learning driven by the 8-to-19 fragment identity shift---is independently testable.

\subsection{Why Candidate-Level Scoring Succeeds Where Query-Level Routing Fails}

Query-level routing strategies consistently failed to improve deployment performance. One conditional-prior gate and a query router (AUC $\approx$ 0.61) both underperformed per-candidate scoring. The predictive signal for fragment replacement is not primarily a property of a query as a whole, but of individual query--candidate interactions. This suggests a testable design principle for related ranking tasks where evidence is distributed across individual items.

\subsection{Prior\_Ranks as Non-Transferable Shortcuts}

Prior\_ranks features encode within-query rank positions whose meaning changes with fragment composition (8 dev $\rightarrow$ 19 blind). The retained model-rank family (F3) contributes positively because base-ranker ranks summarize validated model outputs with stable ranking behavior across fragment identities, whereas prior\_ranks derive from sparse prior-score transforms tightly coupled to specific fragment types. The design rule: prefer raw scores, rates, and intrinsic molecular properties over sparse prior-score rank transforms under fragment distribution shift.

\subsection{Dual-Mode Workflow and Computational Review}

The provenance-label design (G2/G3/G4) provides structured language for communicating the reliability-novelty trade-off. Three scope limitations: (1) rule-based filters only; (2) no experimental validation; (3) fragment-vocabulary-specific.

\subsection{Relationship to Generative and Database-Driven Approaches}

Our closed-vocabulary ranking formulation occupies a distinct position from database-driven tools (SwissBioisostere, NeBULA) and generative approaches (DeepBioisostere, GraphBioisostere). Extending to open-vocabulary ranking with graph-based representations is a natural direction.

\subsection{Limitations and Future Work}

Six limitations: (1) single data source (ChEMBL37K); (2) structure-derived labels, no activity preservation; (3) transform-heldout does not eliminate all chemical relatedness; (4) post-audit feature selection; (5) A4C unvalidated; (6) single-benchmark workflow demonstration. Future directions: independent MMP datasets, activity-based endpoints, GNN integration, prospective pre-registration, expert review calibration.

% ============================================================
%  6. CONCLUSION
% ============================================================

\section{Conclusion}

We introduced a leakage-controlled benchmark and a candidate-level scoring framework for closed-vocabulary scaffold-conditioned fragment replacement. Under the transform-heldout evaluation protocol, a post-audit 77-feature scorer achieves Top-10 = 0.9243 on 13,347 secondary blind queries, $\Delta$Top-10 = +0.0686 over the Score Blend baseline. The central methodological finding is that the most impactful feature-engineering decision is a deletion: removing five per-query prior-rank features eliminates a non-transferable shortcut (+0.0393, 5.9$\times$ hit-loss reduction). Because this deletion was identified after blind diagnostics, the result should be read as a locked post-selection finding that motivates independent prospective replication.

The dual-mode workflow translates these ranking outputs into a practical reporting framework with provenance-stratified alert profiles. Several limitations define the scope: structure-derived labels, single data source, post-audit feature selection, unvalidated A4C proxy. Addressing these limitations through independent replication, activity-based endpoints, and expert review calibration constitutes the natural next phase of this work.

% ============================================================
%  DATA AND SOFTWARE AVAILABILITY
% ============================================================

\section*{Data and Software Availability}

The analysis uses ChEMBL37K, a curated subset of ChEMBL33. A public source-code and evidence archive is available at \url{https://github.com/user141514/paper1/tree/codex/jcim-algorithm-archive}, with the curated experiment-code folder at \texttt{bioisosteric\_diffusion/paper1\_experiment\_code\_archive}. The code-archive commit is \texttt{b750de266d6d47e63de0eeaa945cc999e6f3e08c}; the manuscript/supplement synchronization commit is \texttt{72b7807948c9a1432b99bf1529d98b3d994a3d1f}. The submission-candidate package is frozen under the repository tag \texttt{jcim-submission-candidate-20260601}. An archival DOI will be provided upon acceptance.

% ============================================================
%  REFERENCES
% ============================================================

\begin{thebibliography}{99}



\end{thebibliography}

\end{document}
